%% kapa: exam=1 round=31 year=2020 period=2 form=A subject=appraisal_law subject_ko=감정평가관계법규 seq=21 number=21
%% source: Q-Net 감정평가사 기출문제 (cst00309) artlSeq=5207219 file=extracted/1008/5207219/A형/A형/2020년도 제31회 감정평가사 1차 시험_2교시_A형.pdf
%% conversion: kordoc 4.24.51 -> tools/gen1.py (auto); review=proofread-sonnet; flags=none
\begin{question}{21}
국유재산법령상 중앙관서의 장이 행정재산의 사용료를 면제할 수 있는 경우에 해당하지 않는 것은?
\begin{choices}
  \item 행정재산으로 할 목적으로 기부를 받은 재산에 대하여 기부자의 상속인에게 사용허가하는 경우
  \item 건물 등을 신축하여 기부채납을 하려는 자가 신축기간에 그 부지를 사용하는 경우
  \item 행정재산을 직접 공공용으로 사용하려는 지방자치단체에 사용허가하는 경우
  \item 사용허가를 받은 행정재산을 천재지변으로 사용하지 못하게 되었을 때 그 사용하지 못한 기간에 대한 사용료의 경우
  \item 법령에 따라 정부가 자본금의 50퍼센트 이상을 출자하는 법인이 행정재산을 직접 비영리 공익사업용으로 사용하고자 하여 사용허가하는 경우
\end{choices}
\end{question}
\answer{5}
